
\section{The roadmap focuses on extending ArmoniK capabilities to meet evolving user needs.}

%% ── Frame 1: Roadmap Overview ──
\begin{frame}{Roadmap: Where ArmoniK Is Heading}
  \centering
  \begin{tikzpicture}[
    pillar/.style={
      rectangle, rounded corners=5pt,
      minimum width=4.8cm, minimum height=2.8cm,
      inner sep=8pt, align=left
    },
    convergence/.style={
      rectangle, rounded corners=5pt,
      fill=aneoorange!15, draw=aneoorange!40,
      text=aneoblack, font=\footnotesize\bfseries,
      minimum width=5.5cm, minimum height=0.8cm,
      align=center
    }
  ]
    % Left pillar — Easier to Use
    \node[pillar, fill=aneoturquoise!12, draw=aneoturquoise!40,
          anchor=north east] (left) at (-0.3,0)
      {\parbox{4.4cm}{{\footnotesize\bfseries\color{aneoturquoise!80!black}Easier to Use}\\[4pt]
       {\tiny
        --~Deployment portability\\
        --~Security \& access control\\
        --~Observability \& testing}}};

    % Right pillar — More Performant
    \node[pillar, fill=aneoblue!12, draw=aneoblue!35,
          anchor=north west] (right) at (0.3,0)
      {\parbox{4.4cm}{{\footnotesize\bfseries\color{aneoblue}More Performant}\\[4pt]
       {\tiny
        --~Persistence \& control plane\\
        --~Scaling \& multi-region\\
        --~Scheduling \& data management}}};

    % Convergence box
    \node[convergence] (conv) at (0,-3.8)
      {Managed PaaS\;\;\textperiodcentered\;\;Custom Deployable};

    % Arrows from pillars to convergence
    \draw[-Stealth, thick, aneoblack!40] (left.south)  -- (conv.north -| left.south);
    \draw[-Stealth, thick, aneoblack!40] (right.south) -- (conv.north -| right.south);
  \end{tikzpicture}

  \note{
    \begin{itemize}
      \item The roadmap is organized around two pillars: making ArmoniK easier
            to use and making it more performant. These aren't separate tracks —
            they reinforce each other. A faster platform that's hard to deploy
            doesn't help anyone, and an easy platform that can't perform doesn't
            meet our users' needs.
      \item The dual delivery model is the strategic target that ties everything
            together. We're building ArmoniK so it can be consumed as a turnkey
            managed PaaS — where we operate the control plane and you keep your
            data and compute in your own tenant — while also remaining fully
            deployable in custom environments: on-premises, bare metal,
            supercomputers. Every roadmap topic contributes to one or both of
            these delivery models.
      \item Let me walk you through the five thematic groups. We'll start with
            the core engine — persistence and control plane — then move outward
            to scaling, deployment portability, security and observability, and
            finally advanced scheduling and ecosystem integration.
    \end{itemize}
  }
\end{frame}

%% ── Frame 2: Persistence, Control Plane & Performance Optimization ──
\begin{frame}{Core Engine: Persistence \& Control Plane}
  \centering
  \begin{tikzpicture}[
    layer/.style={
      rectangle, rounded corners=4pt,
      text width=10.5cm, align=left,
      inner sep=6pt, minimum height=1.0cm
    }
  ]
    % Bottom layer — SQL Migration (darkest)
    \node[layer, fill=aneoblue, text=white] (sql) at (0,0)
      {{\footnotesize\bfseries SQL Migration}\\[2pt]
       {\tiny From NoSQL to SQL\;\textperiodcentered\;Consistency \& scalability\;\textperiodcentered\;Managed DB services}};

    % Layer 2 — Control Plane Refactoring
    \node[layer, fill=aneoblue!25, draw=aneoblue!40, text=aneoblack,
          anchor=south] (refactor) at ([yshift=4pt]sql.north)
      {{\footnotesize\bfseries Control Plane Refactoring}\\[2pt]
       {\tiny Stateless scheduling\;\textperiodcentered\;Decoupled from DB \& queue\;\textperiodcentered\;Idempotent RPC APIs}};

    % Layer 3 — Control Plane Optimization
    \node[layer, fill=aneoblue!15, draw=aneoblue!30, text=aneoblack,
          anchor=south] (optim) at ([yshift=4pt]refactor.north)
      {{\footnotesize\bfseries Control Plane Optimization}\\[2pt]
       {\tiny Batched operations\;\textperiodcentered\;Task prefetching\;\textperiodcentered\;Reduced database round-trips}};

    % Top layer — Data Plane Isolation (lightest)
    \node[layer, fill=aneoblue!8, draw=aneoblue!20, text=aneoblack,
          anchor=south] (dataplane) at ([yshift=4pt]optim.north)
      {{\footnotesize\bfseries Data Plane Isolation}\\[2pt]
       {\tiny Dedicated Data Agent\;\textperiodcentered\;Separation of control and data traffic}};

    % Left-side rotated annotation
    \node[font=\tiny\itshape, text=aneoblack!50, rotate=90]
      at ([xshift=-6.2cm]{{barycentric cs:sql=0.5,dataplane=0.5}})
      {Foundation\;$\rightarrow$\;Surface};
  \end{tikzpicture}

  \note{
    \begin{itemize}
      \item The foundation of everything is the SQL migration. We're moving from
            MongoDB to a SQL relational database. Why? Three reasons: strong
            consistency instead of eventual consistency, better scalability
            patterns for our access patterns, and --- critically --- easier access
            to managed database services. Whether you're in the cloud using RDS
            or Aurora, or on-premises where your IT department already operates
            PostgreSQL, SQL is the lingua franca. This directly enables the
            managed PaaS model.
      \item On top of that, we're refactoring the control plane for horizontal
            scalability. Today, agents are coupled to the database and queue.
            We're decoupling them into a stateless scheduling service with
            idempotent RPC APIs. This means you can scale the control plane
            horizontally --- add more instances under load --- without worrying
            about state conflicts.
      \item The optimization layer is about raw throughput: batching database
            operations, prefetching tasks, reducing the number of database
            round-trips per scheduling cycle. And at the top, data plane
            isolation gives us a dedicated Data Agent that separates control
            traffic from data traffic, preventing large data transfers from
            starving the scheduler.
    \end{itemize}
  }
\end{frame}

%% ── Frame 3: Cloud-Native Scaling, Multi-Region & Disaster Recovery ──
\begin{frame}{Scaling \& Multi-Region Resilience}
  \centering
  \begin{tikzpicture}[
    regionbox/.style={
      rectangle, rounded corners=6pt,
      draw=aneoturquoise!50, fill=aneoturquoise!5,
      minimum width=4.5cm, minimum height=3.5cm
    },
    regionlabel/.style={
      font=\footnotesize\bfseries, text=aneoturquoise!80!black
    },
    partbox/.style={
      rectangle, rounded corners=3pt,
      fill=aneoturquoise!20, draw=aneoturquoise!40,
      font=\tiny\bfseries, text=aneoblack,
      minimum width=3.2cm, minimum height=0.6cm, align=center
    },
    storagebox/.style={
      rectangle, rounded corners=3pt,
      fill=aneoblue!10, draw=aneoblue!30,
      font=\tiny, text=aneoblack,
      minimum width=3.2cm, minimum height=0.5cm, align=center
    },
    routingnode/.style={
      rectangle, rounded corners=4pt,
      fill=aneoorange!20, draw=aneoorange!50,
      font=\tiny\bfseries, text=aneoblack,
      minimum width=2.0cm, minimum height=0.7cm, align=center
    }
  ]
    % Region A (left)
    \node[regionbox] (regA) at (-4,0) {};
    \node[regionlabel, anchor=north] at (regA.north) {Region A};
    \node[partbox]    (partA) at ([yshift=0.2cm]regA.center) {Partitions + Autoscaling};
    \node[storagebox] (storA) at ([yshift=-1.0cm]regA.center) {SQL + Obj Storage};

    % Region B (right)
    \node[regionbox] (regB) at (4,0) {};
    \node[regionlabel, anchor=north] at (regB.north) {Region B};
    \node[partbox]    (partB) at ([yshift=0.2cm]regB.center) {Partitions + Autoscaling};
    \node[storagebox] (storB) at ([yshift=-1.0cm]regB.center) {SQL + Obj Storage};

    % Central routing node
    \node[routingnode] (routing) at (0,0.2) {Workload\\Routing};

    % Bidirectional arrows: routing ↔ regions
    \draw[-Stealth, thick, aneoorange!70] (routing.west) -- (partA.east);
    \draw[-Stealth, thick, aneoorange!70] (partA.east)   -- (routing.west);
    \draw[-Stealth, thick, aneoorange!70] (routing.east) -- (partB.west);
    \draw[-Stealth, thick, aneoorange!70] (partB.west)   -- (routing.east);

    % Dashed replication arrow between storage boxes
    \draw[-Stealth, dashed, aneoblue!50] (storA.east) -- (storB.west)
      node[midway, above, font=\tiny\itshape, text=aneoblue!60] {replication};

    % Region Bursting annotation
    \node[font=\tiny\itshape, text=aneoorange!80] at (0,-2.6)
      {Region Bursting\quad{\upshape(isolated data planes)}};
  \end{tikzpicture}

  \note{
    \begin{itemize}
      \item This frame is about making ArmoniK a globally resilient platform.
            It starts with a Kubernetes-native autoscaling operator that uses
            Prometheus metrics and Karpenter to scale across multiple clusters
            and regions. Combined with automated partition provisioning — where
            partitions are declared and a partition operator handles the rest —
            this gives you elastic, self-managing infrastructure.
      \item The routing layer is where it gets interesting. Workload routing
            uses metrics to decide where sessions land — which region, which
            cluster. It supports region-aware placement for latency and
            compliance, and on-prem to cloud overflow when local capacity is
            exhausted. This is critical for hybrid deployments.
      \item Disaster recovery is built on three pillars: globally distributed
            SQL for metadata, rebuildable queues that can be reconstructed from
            the database, and cross-region object storage replication for data.
            Region bursting extends this further — when one region is saturated,
            workloads overflow to another region with isolated data planes, so
            there's no data leakage between regions.
    \end{itemize}
  }
\end{frame}

%% ── Frame 4: Deployment Portability ──
\begin{frame}{Deployment Portability}
  \centering
  \begin{tikzpicture}[
    card/.style={
      rectangle, rounded corners=5pt,
      minimum width=2.6cm, minimum height=2.4cm,
      text width=2.2cm, align=center, inner sep=5pt
    }
  ]
    % Spectrum line with arrowheads
    \draw[draw=aneoblack!30, thick, {Stealth}-{Stealth}]
      (-6,0) -- (6,0);
    \node[font=\tiny\bfseries, text=aneoblack!60, anchor=east]
      at (-6.15,0) {Managed};
    \node[font=\tiny\bfseries, text=aneoblack!60, anchor=west]
      at (6.15,0) {Custom};

    % Card 1 — Cloud Marketplace
    \node[card, fill=aneoorange!20, draw=aneoorange!50] (mp) at (-4.2,-2.0)
      {{\footnotesize\bfseries\color{aneoblack}Cloud\\Marketplace}\\[4pt]
       {\tiny Billing integration\;\textperiodcentered\;Simplified adoption}};

    % Card 2 — Cloud IaaS / K8s
    \node[card, fill=aneoorange!12, draw=aneoorange!35] (iaas) at (-1.4,-2.0)
      {{\footnotesize\bfseries\color{aneoblack}Cloud\\IaaS / K8s}\\[4pt]
       {\tiny Helm + Terraform\;\textperiodcentered\;Reference deployments}};

    % Card 3 — Bare Metal
    \node[card, fill=aneoorange!8, draw=aneoorange!25] (bare) at (1.4,-2.0)
      {{\footnotesize\bfseries\color{aneoblack}Bare Metal}\\[4pt]
       {\tiny Process-based lifecycle\;\textperiodcentered\;No Kubernetes required}};

    % Card 4 — Supercomputer
    \node[card, fill=aneoorange!5, draw=aneoorange!15] (hpc) at (4.2,-2.0)
      {{\footnotesize\bfseries\color{aneoblack}Supercomputer}\\[4pt]
       {\tiny Slurm integration\;\textperiodcentered\;HPC cluster environments}};

    % Connector ticks from spectrum line to cards
    \draw[aneoblack!20] (-4.2,0) -- (-4.2,-0.8);
    \draw[aneoblack!20] (-1.4,0) -- (-1.4,-0.8);
    \draw[aneoblack!20] (1.4,0)  -- (1.4,-0.8);
    \draw[aneoblack!20] (4.2,0)  -- (4.2,-0.8);

    % Dual Delivery bracket — simple line approach (no decorations library)
    % Managed PaaS span
    \draw[thick, aneoturquoise!60] (-5.5,-3.7) -- (-5.5,-3.9) -- (-0.1,-3.9) -- (-0.1,-3.7);
    \node[font=\tiny\bfseries, text=aneoturquoise!80!black] at (-2.8,-4.25)
      {Managed PaaS};

    % Custom Deployable span
    \draw[thick, aneoblue!50] (0.1,-3.7) -- (0.1,-3.9) -- (5.5,-3.9) -- (5.5,-3.7);
    \node[font=\tiny\bfseries, text=aneoblue] at (2.8,-4.25)
      {Custom Deployable};
  \end{tikzpicture}

  \note{
    \begin{itemize}
      \item Deployment portability is what makes the dual delivery model real.
            On the left end of the spectrum, we have the cloud marketplace —
            one-click deployment with billing integration, the fastest path to
            getting started. Next to it, the standard cloud IaaS deployment
            with Helm and Terraform that our current users know well.
      \item Moving right, bare metal deployment is about running ArmoniK
            without Kubernetes at all. Some environments — especially in
            regulated industries or legacy infrastructure — can't or won't
            run Kubernetes. We're building a process-based lifecycle manager
            that handles worker provisioning and health monitoring directly
            on bare metal.
      \item At the far right, supercomputer deployment integrates with Slurm,
            the dominant job scheduler in HPC. This opens ArmoniK to the
            scientific computing community — nuclear fusion simulations,
            seismic imaging, automotive crash simulations — where
            supercomputers are the primary compute resource. Together, these
            four deployment models mean ArmoniK truly runs anywhere.
    \end{itemize}
  }
\end{frame}

%% ── Frame 5: Security, Quality & Observability ──
\begin{frame}{Security, Quality \& Observability}
  \begin{columns}[T]

    %% Left column — Security & Access Control
    \begin{column}{0.48\textwidth}
      \centering
      \begin{tikzpicture}[
        card/.style={
          rectangle, rounded corners=4pt,
          text width=4.4cm, minimum width=4.8cm,
          minimum height=1.1cm, align=left, inner sep=5pt
        }
      ]
        % Card 1 — Secure Data Access
        \node[card, fill=aneoblue!12, draw=aneoblue!35] (sda) at (0,0)
          {{\footnotesize\bfseries\color{aneoblue}Secure Data Access}\\[2pt]
           {\tiny Token-based authorization\;\textperiodcentered\;Tenant isolation\;\textperiodcentered\;Decoupled data management}};

        % Card 2 — User Management & Access Control
        \node[card, fill=aneoblue!8, draw=aneoblue!25,
              anchor=north] (umac) at ([yshift=-6pt]sda.south)
          {{\footnotesize\bfseries\color{aneoblue!70!black}User Management \& Access Control}\\[2pt]
           {\tiny RBAC/ABAC\;\textperiodcentered\;SSO \& identity federation\;\textperiodcentered\;Multi-tenant isolation}};
      \end{tikzpicture}
    \end{column}

    %% Right column — Quality & Observability
    \begin{column}{0.48\textwidth}
      \centering
      \begin{tikzpicture}[
        card/.style={
          rectangle, rounded corners=4pt,
          text width=4.4cm, minimum width=4.8cm,
          minimum height=1.1cm, align=left, inner sep=5pt
        }
      ]
        % Card 3 — Non-Regression Testing
        \node[card, fill=aneoturquoise!10, draw=aneoturquoise!30] (nrt) at (0,0)
          {{\footnotesize\bfseries\color{aneoturquoise!80!black}Non-Regression Testing}\\[2pt]
           {\tiny Automated test suites\;\textperiodcentered\;Performance benchmarks\;\textperiodcentered\;Release quality gates}};

        % Card 4 — Functional Observability
        \node[card, fill=aneoturquoise!15, draw=aneoturquoise!40,
              anchor=north] (fobs) at ([yshift=-6pt]nrt.south)
          {{\footnotesize\bfseries\color{aneoturquoise!80!black}Functional Observability}\\[2pt]
           {\tiny Critical path analysis\;\textperiodcentered\;Task metrics comparison across runs}};

        % Card 5 — Unified Infrastructure Observability
        \node[card, fill=aneoturquoise!20, draw=aneoturquoise!50,
              anchor=north] (iobs) at ([yshift=-6pt]fobs.south)
          {{\footnotesize\bfseries\color{aneoturquoise!80!black}Unified Infrastructure Observability}\\[2pt]
           {\tiny Centralized dashboard\;\textperiodcentered\;Hardware, network \& application metrics}};
      \end{tikzpicture}
    \end{column}

  \end{columns}

  \note{
    \begin{itemize}
      \item Security is non-negotiable for our users in regulated industries.
            The secure data access architecture introduces token-based
            authorization with tenant isolation — each tenant's data is
            cryptographically separated. User management adds RBAC and ABAC
            with SSO integration, so enterprises can plug ArmoniK into their
            existing identity infrastructure.
      \item Non-regression testing is about operational confidence. Every
            release goes through automated test suites and performance
            benchmarks. If a change degrades throughput or introduces a
            regression, we catch it before it reaches production. This is
            especially important for the managed PaaS model where we're
            responsible for the control plane.
      \item Observability splits into two layers. Functional observability
            gives you insight into your algorithms — critical path analysis,
            comparing task metrics across runs to spot performance regressions
            in your code. Infrastructure observability gives operators a
            centralized dashboard covering hardware utilization, network
            health, and application metrics. Together, they close the loop
            from "something is slow" to "here's exactly why."
    \end{itemize}
  }
\end{frame}

%% ── Frame 6: Advanced Scheduling, Data Management & Ecosystem Integration ──
\begin{frame}{Scheduling, Data Management \& Ecosystem}
  \begin{columns}[T]

    %% Left column — Scheduling & Data Management
    \begin{column}{0.55\textwidth}
      \centering
      \begin{tikzpicture}[
        card/.style={
          rectangle, rounded corners=4pt,
          text width=5.0cm, align=left, inner sep=5pt,
          minimum width=5.4cm, minimum height=1.0cm
        }
      ]
        % Card 1 — Advanced Scheduling
        \node[card, fill=aneoturquoise!15, draw=aneoturquoise!40] (sched) at (0,0)
          {{\footnotesize\bfseries\color{aneoturquoise!80!black}Advanced Scheduling}\\[2pt]
           {\tiny\color{aneoblack}Fair share\;\textperiodcentered\;Metric-driven\;\textperiodcentered\;Task \& data affinity}};

        % Card 2 — Data Dependencies
        \node[card, fill=aneoturquoise!10, draw=aneoturquoise!30,
              anchor=north] (deps) at ([yshift=-10pt]sched.south)
          {{\footnotesize\bfseries\color{aneoturquoise!80!black}Data Dependencies}\\[2pt]
           {\tiny\color{aneoblack}Optional\;\textperiodcentered\;Lazy\;\textperiodcentered\;Lazy-optional categories}};

        % Card 3 — Task Checkpointing
        \node[card, fill=aneoturquoise!8, draw=aneoturquoise!25,
              anchor=north] (ckpt) at ([yshift=-10pt]deps.south)
          {{\footnotesize\bfseries\color{aneoturquoise!80!black}Task Checkpointing}\\[2pt]
           {\tiny\color{aneoblack}Save \& restore state for fault recovery}};

        % Connecting arrows
        \draw[-Stealth, thick, aneoturquoise!60] (sched.south) -- (deps.north);
        \draw[-Stealth, thick, aneoturquoise!60] (deps.south)  -- (ckpt.north);
      \end{tikzpicture}
    \end{column}

    %% Right column — OpenGRIS Integration
    \begin{column}{0.42\textwidth}
      \centering
      \begin{tikzpicture}
        \node[rectangle, rounded corners=5pt,
              fill=aneoorange!15, draw=aneoorange!40,
              text width=4.0cm, align=left, inner sep=8pt,
              minimum width=4.4cm, minimum height=3.0cm] (ogris) at (0,0)
          {{\footnotesize\bfseries\color{aneoorange!80!black}OpenGRIS Integration}\\[6pt]
           {\small
            \textbullet\;ParFun SDK\\
            \textbullet\;ParGraph SDK\\
            \textbullet\;Scaler Metascheduler}\\[6pt]
           {\tiny\itshape\color{aneoblack!50}Ecosystem reach}};
      \end{tikzpicture}
    \end{column}

  \end{columns}

  \note{
    \begin{itemize}
      \item Advanced scheduling is about making ArmoniK smarter about where
            and when tasks run. Fair share scheduling ensures no single user
            monopolizes the cluster. Metric-driven scheduling uses runtime
            telemetry to make placement decisions. And task-data affinity
            scheduling co-locates tasks with their data — this ties directly
            back to the data-aware architecture we discussed in the vision
            section.
      \item Data dependency management introduces three new categories beyond
            the current mandatory dependencies: optional data that tasks can
            proceed without, lazy data that's fetched on demand rather than
            prefetched, and lazy-optional combining both. This gives developers
            fine-grained control over the performance-correctness tradeoff.
            Task checkpointing complements this by letting long-running tasks
            save their state periodically, so a failure doesn't mean starting
            from scratch.
      \item OpenGRIS integration is our bridge to the broader HPC ecosystem.
            ParFun and ParGraph are SDKs for expressing parallel computations,
            and the Scaler metascheduler enables ArmoniK to participate in
            multi-runtime scheduling. This positions ArmoniK not as an
            isolated platform but as a first-class citizen in the emerging
            HPC standards landscape.
    \end{itemize}
  }
\end{frame}

%% ── Frame 7: Roadmap Summary ──
\begin{frame}{Roadmap Summary}
  \centering
  \begin{tikzpicture}[
    banner/.style={
      rectangle, rounded corners=5pt,
      fill=aneoblue, text=white,
      font=\small\bfseries,
      text width=10cm, align=center,
      minimum height=0.7cm, inner sep=5pt
    },
    grouplabel/.style={
      font=\tiny\bfseries, text=aneoblack
    },
    dualbox/.style={
      rectangle, rounded corners=5pt,
      minimum width=5.5cm, minimum height=1.6cm,
      draw=aneoblack!30, inner sep=0pt
    }
  ]
    % Top banner
    \node[banner] (banner) at (0,0)
      {Easier to use\;\textperiodcentered\;More performant};

    % Five thematic group labels on the left
    \node[grouplabel] (g1) at (-3.8,-1.4)
      {\tikz\fill[aneoblue] (0,0) circle (2.5pt); \;Persistence \& Control Plane};
    \node[grouplabel] (g2) at (-3.8,-1.9)
      {\tikz\fill[aneoturquoise] (0,0) circle (2.5pt); \;Scaling \& Multi-Region};
    \node[grouplabel] (g3) at (-3.8,-2.4)
      {\tikz\fill[aneoorange] (0,0) circle (2.5pt); \;Deployment Portability};
    \node[grouplabel] (g4) at (-3.8,-2.9)
      {\tikz\fill[aneoblue!70!black] (0,0) circle (2.5pt); \;Security \& Observability};
    \node[grouplabel] (g5) at (-3.8,-3.4)
      {\tikz\fill[aneoturquoise!80!black] (0,0) circle (2.5pt); \;Scheduling \& Ecosystem};

    % Dual Delivery box (outer container)
    \node[dualbox] (dual) at (2.8,-2.4) {};

    % Top half — Managed PaaS
    \node[rectangle, rounded corners=5pt,
          fill=aneoturquoise!20, text=aneoturquoise!80!black,
          font=\footnotesize\bfseries,
          minimum width=5.3cm, minimum height=0.7cm,
          anchor=south, inner sep=4pt]
      (paas) at ([yshift=1pt]dual.center) {Managed PaaS};

    % Bottom half — Custom Deployable
    \node[rectangle, rounded corners=5pt,
          fill=aneoorange!15, text=aneoorange!80!black,
          font=\footnotesize\bfseries,
          minimum width=5.3cm, minimum height=0.7cm,
          anchor=north, inner sep=4pt]
      (custom) at ([yshift=-1pt]dual.center) {Custom Deployable};

    % Convergence arrows from labels to Dual Delivery box
    \draw[-Stealth, thick, aneoblack!40] (g1.east) -- (dual.west |- g1.east);
    \draw[-Stealth, thick, aneoblack!40] (g2.east) -- (dual.west |- g2.east);
    \draw[-Stealth, thick, aneoblack!40] (g3.east) -- (dual.west |- g3.east);
    \draw[-Stealth, thick, aneoblack!40] (g4.east) -- (dual.west |- g4.east);
    \draw[-Stealth, thick, aneoblack!40] (g5.east) -- (dual.west |- g5.east);

    % Bottom annotations
    \node[font=\small\itshape, text=aneoblack!60, align=center]
      at (0,-4.6)
      {A vision and priority list, not a timeline\\[2pt]
       Customer collaboration accelerates specific topics};
  \end{tikzpicture}

  \note{
    \begin{itemize}
      \item Let me bring it all together. Everything we've discussed --- the SQL
            migration, the control plane refactoring, multi-region resilience,
            deployment portability, security, observability, advanced scheduling
            --- all of it converges on one strategic target: delivering ArmoniK
            both as a managed PaaS and as a custom deployable solution.
      \item I want to be clear about what this roadmap is and isn't. It's a
            vision and a priority list. It's not a Gantt chart with fixed
            dates. Development moves forward continuously --- we ship every
            4 to 6 weeks --- and we onboard new customers at the same time.
            Priorities shift based on what our users need most.
      \item And that's the last point: customer collaboration shapes this
            roadmap. Some topics are accelerated because a specific customer
            needs them. If you see something on this roadmap that matters to
            your organization, talk to us --- that's how topics move up the
            priority list. With that, let's open it up for questions.
    \end{itemize}
  }
\end{frame}
